\documentclass[conference]{IEEEtran}
\IEEEoverridecommandlockouts

% Archivo fuente en español del manuscrito de la tesis/artículo.
% El contenido de trabajo actual se mantiene aquí y se sincroniza con Overleaf.

\usepackage{cite}
\usepackage{amsmath,amssymb,amsfonts}
\usepackage{algorithmic}
\usepackage{graphicx}
\usepackage{textcomp}
\usepackage{xcolor}
\usepackage{float}
\usepackage{booktabs}
\usepackage{array}
\usepackage{multirow}
\usepackage{url}
\usepackage{siunitx}
\usepackage{balance}

\def\BibTeX{{\rm B\kern-.05em{\sc i\kern-.025em b}\kern-.08em
    T\kern-.1667em\lower.7ex\hbox{E}\kern-.125emX}}

\begin{document}

\title{Diseño y desarrollo de un banco de pruebas multifuncional para la caracterización térmica y vibratoria de sensores capacitivos industriales y la evaluación de un adaptador termomecánico con elemento SMA}

% IMPORTANTE: aquí se conserva el manuscrito de trabajo.
% Copiar/actualizar este archivo con la versión completa de Main_es.tex.

\author{
\IEEEauthorblockN{F.E. Rabanal Medina}
\IEEEauthorblockA{\textit{Universidad Peruana de Ciencias Aplicadas}\\
U201920529@upc.edu.pe\\
Lima, Perú}
}

\maketitle

\begin{abstract}
Versión de trabajo. El manuscrito completo se incorpora en la rama de desarrollo documental y se mantiene sincronizado con el proyecto de Overleaf.
\end{abstract}

\bibliographystyle{IEEEtran}
\bibliography{../../references/references}

\end{document}
